
\subsection{Hardware-Aware Efficiency Metrics}

Compression is evaluated through several distinct efficiency metrics. Parameter count, memory footprint, arithmetic cost, latency, and throughput capture different aspects of model execution. Hardware-aware evaluation therefore requires separate treatment of storage cost, arithmetic cost, and runtime behavior~\cite{Sze2017EfficientDNN}.

\subsubsection{Parameter Count and Model Size}

Let $\params = \{W^{(l)}, b^{(l)}\}_{l=1}^{L}$ denote the parameters of a network. The total parameter count is definded by Eq.\eqref{eq:parameter_count}:
\begin{equation}
	\label{eq:parameter_count}
	P = \sum_{l=1}^{L} \left( |W^{(l)}| + |b^{(l)}| \right).
\end{equation}
If the parameters are stored with precision $b_w$ bits. Eq.\eqref{eq:model_size} gives the model size in bytes:
\begin{equation}
	\label{eq:model_size}
	M_{\text{model}} = \frac{P\, b_w}{8} \quad \text{bytes}.
\end{equation}
This follows from counting the stored elements and multiplying by the number of bytes used by each element~\cite{PyTorchNumel,PyTorchElementSize}.
This quantity determines whether a model fits within device memory and how much parameter traffic must be sustained during inference. Quantization reduces $b_w$, pruning reduces the number of active parameters, and low-rank factorization reduces the parameterization itself.

\subsubsection{MACs, FLOPs, and Arithmetic Intensity}

For a dense matrix multiplication $Y = AB$ with $A \in \reals^{m \times k}$ and $B \in \reals^{k \times n}$, the arithmetic cost is defined by Eq.~\eqref{eq:arithmetic_cost}:
\begin{equation}
	\label{eq:arithmetic_cost}
	\mathrm{MACs}(A,B) = mkn,
	\qquad
	\mathrm{FLOPs}(A,B) \approx 2mkn.
\end{equation}

One multiply-accumulate corresponds to one MAC, while FLOP counting often treats the multiply and the addition as two floating-point operations~\cite{NVIDIAMatrixMultiplication}. In Transformer layers, the feed-forward subnetwork contributes large dense matrix products, while self-attention adds the sequence-dependent $O(n^2 d_{\text{model}})$ cost of score computation and value aggregation.


\subsubsection{Latency, and Throughput}


Latency and throughput characterize runtime from complementary viewpoints. Latency is the time required to complete one inference request or one decoding step~\cite{Sze2017EfficientDNN} as shown in Eq.\eqref{eq:latency}:
\begin{equation}
	\label{eq:latency}
	T_{\text{latency}} = T_{\text{compute}} + T_{\text{memory}} + T_{\text{overhead}}.
\end{equation}
Throughput measures completed work per unit time, for example Eq.\eqref{eq:throughput} gives
requests per second or generated tokens per second~\cite{Sze2017EfficientDNN}:
\begin{equation}
	\label{eq:throughput}
	\mathrm{Throughput} = \frac{N_{\text{outputs}}}{T}.
\end{equation}
